\introduction{The Measure of All Structures}
Puella Magi Madoka Magica (hereafter generally Madoka Magica to distinguish it from the title character, Madoka Kaname \cite{ref1}, who shall generally be referred to as simply Madoka) is many things. It is, most obviously, an anime series of the magical girl genre by Studio Shaft, but beyond that it is a critique of and homage to that genre; it is an attempt at self-redemption by a writer lost in the dark; it is a Buddhist retelling of Faust. It is a study in contrasts, at once a feminist examination of questions of consent, privilege, and gender roles, and a sexist objectification of the underage female form. It is about sacrifice and selfishness, despair and hope, love and death, decay and life, and how all of these seeming opposites are really one and the same.

Most of all, though, it is about people, who make choices (frequently bad ones), and suffer, and strive, and succeed, and fail. It is about characters who love each other and hate each other and don't care about each other, sometimes all at once.

Over the course of this book, I will be talking quite a bit about grand themes, philosophies, and symbols, so here at the beginning I wanted to establish a reminder: This is a story about people, because all stories are about people \cite{ref2}. 
\paragraph{Why Madoka?}~{}

In a mere twelve 22-minute episodes and one feature film (there were three released, but the first two are arguably just compilations of the TV series with a few minor tweaks), Madoka Magica packs in an enormous amount of story, with rich characters, complex themes, and moments of profound sadness, apprehension, anger, and joy. Analyzing it is tremendously rewarding, as it reveals more layers every time it is examined. In particular, as very much a postmodern work itself, it greatly rewards analysis with a postmodern bent.
\paragraph{Po-wha huh?}~{}

Postmodernism is a rather complex term to define, but I shall essay that task, at least to explain what I mean when I talk about it. Philosophically, as the name implies, postmodernism is a step beyond modernism.

In a nutshell, the core realization of modernism is that symbols are fundamentally arbitrary. That is, there is no relationship between the signifier (the thing that does the symbolizing) and the signified (the thing that is symbolized), except in the mind of the person looking at the signifier. So, for example, a red hexagon means "stop" not because there is some logical connection between redness, hexagonality, and stopping, but because somebody somewhere decided that red hexagons should mean stop, and convinced others to go along with it. This is the essential concept of social constructionism, that the concepts and symbols with which we understand our world are constructed from social circumstances and relationships. It naturally follows that these symbols are therefore as fluid as societies and relationships are, and do not have fixed, objective meaning \cite{ref3}. 

Since all art (and language) is a series of symbols according to postmodernist thought, it follows that the meaning of any given work is in some sense arbitrary, a product of social construction rather than a logically necessary relation. Traditionally, art got around this by using agreed-upon, culturally defined symbols, such as words or representational images (using the image of an apple to stand in for an apple, for example). These act as guide rails of a sort, allowing the person experiencing the art to start with a few familiar symbols, then build from there. Modernism being largely characterized by a rejection of tradition and a sense of disillusionment and disintegration, modern art often dispenses with some or all of these guide rails, aggressively challenging the very idea that art has—or should have—non-arbitrary meaning. Examples include James Joyce's works, presenting stream-of-consciousness text without normal sentence structure, and Mondrian's paintings depicting abstract colors and shapes rather than representations of familiar objects. Paradoxically, however, modernism frequently seeks to recreate the order for which it feels that loss \cite{ref4}, as Joyce famously depicts an ordinary day as an epic struggle, or Mondrian arranges his abstract colors into grid-like patterns.

Postmodernism, by contrast, rejects the idea that socially constructed meaning is the same thing as "no meaning." Modernism says "The relationship between signifier and signified is arbitrary, so nothing means anything"; postmodernism says "The relationship between signifier and signified is arbitrary, so we're free to decide what everything means." Thus postmodernist works do not need to reintroduce lost order, since it never existed to begin with. For much the same reason, the postmodern period has seen an embrace of pop culture as a medium for artistic expression \cite{ref5}. Modernism rejects traditional meanings, forms, and techniques in order to create a new order to replace the old; postmodernism combines the traditional and the new in order to pursue the odd or interesting \cite{ref6}. 

Where modern art removes the guide rails, postmodern art tends to function by drawing attention to the guide rails. This is done primarily via a process of decontextualization and re-contextualization, removing familiar signifiers from their usual contexts and placing them within other contexts, for example by mixing elements from multiple genres, creating pastiches of familiar works, or breaking the fourth wall \cite{ref7}. The resulting sense of disorientation is generally then exploited either for humorous or horrific effect, or to provoke thought. (It can, of course, also be funny, frightening, and thought-provoking all at once.)

Postmodernism can be divided into "positive" strands that emphasize playfulness and freedom, and "negative" strands that emphasize disorder and cynicism or despair \cite{ref8}. Madoka Magica falls largely into the latter camp, and makes heavy use of postmodern techniques as part of telling its story, for example (as we'll discuss more in the next chapter) by having the witches take the form of alien, invasive art styles that overwrite the familiar anime style of the show.
\paragraph{This Book}~{}

This book is divided into three parts. The first, "The Television Series," contains twelve chapters, each covering one episode of the television series Puella Magi Madoka Magica, and each titled with a fragment of a quote from Czech playwright, activist, and politician Václav Havel, followed by the title of the episode in parentheses. I have chosen Havel quotes for reasons that should become clear as we progress through the series. The title of this book is also from him: "Perhaps hopelessness is the very soil that nourishes human hope; perhaps one could never find sense in life without first experiencing its absurdity." \cite{ref9} 

Following this first part is an interlude, "One Perfect Moment," on the series' use of narrative elements of Goethe's Faust.

The second and shortest section, "The Comics," contains three chapters, each covering one of the manga spinoffs published during and after the broadcast of the television series. Each of these chapters is titled with a quote from the television series.

Following this part is a second interlude, "Corpse of Milk," on the significance of cheese and the alchemical concept of putrefaction in Puella Magi Madoka Magica the Movie: Rebellion (hereafter generally referred to as Rebellion).

The final section, "Rebellion," consists of seven chapters covering Rebellion. I have chosen to cover only this film, the third in the larger Madoka Magica franchise, as the first two are simply recaps of the television series, with little new footage, let alone story. Rebellion, by contrast, is a sequel to the television series \cite{ref10}. Each of the seven chapters discusses a different possible target for the titular rebellion, and so all their titles are of the form "Against \underline{\makebox[1cm]{}}."

This book is an analytical work, not a guidebook. Plot, character, and production details are described where relevant to the analysis, but generally speaking this book assumes familiarity with the work under discussion. At time of writing, the television series is readily available in the U.S. on DVD, BluRay, and via the Netflix and CrunchyRoll streaming services; the comics discussed have likewise either been published, or are currently being published, in English translation in the U.S. Rebellion is rather more difficult to get ahold of following a brief and limited theatrical run; however, it is possible to import the Japanese BluRay, which has English subtitles, and a BluRay release with both English subtitle and dub options is due for wide release in the U.S. in April 2015. Given this assumption of familiarity, the works under discussion are not cited, as they are the topic of study. Citations are included as needed for references to works outside Puella Magi Madoka Magica and the spinoffs discussed in this book.

Finally, several chapters in this book, most notably Chapters 2 and 9, delve into a feminist analysis of aspects of the series in particular and the magical girl genre in general. Throughout this discussion, please keep in mind that the feminist theory employed is Western in origin and derives from European cultural history and assumptions. As such, while it is applicable to the viewing experience of a Western audience, it is not necessarily applicable to the Japanese culture, which is both the originating culture of the series and its primary audience. The discussion in those chapters should thus be assumed also to apply to the series as perceived from a Western perspective, and is not necessarily generalizable to the Japanese experience.

\paragraph{Acknowledgments}~{}

This book would not have been possible without a number of people. First and foremost is Viga Gadson, my dear friend and cover designer, who introduced me to Madoka Magica in the first place. Second is Unnoun, whose relentless linking of my blog entries on the TV series to various discussion fora led to those entries having among the highest traffic of anything I'd ever written, in turn leading me to write more on Madoka Magica and ultimately this book. Third is my editor, Kit Paige, to whom I owe whatever veneer of respectable scholarship this book may have. Any errors remaining can be attributed to my stubbornness, not her efforts. Fourth are Mark Oshiro of MarkWatches.net and Chuck Sonnenburg of SFDebris.com, whose (very different) reviews of the series' episodes helped me get some much-needed critical distance from my own ideas. Last but not least are the readers of my blog, JedABlue.com, whose insightful comments on the first drafts of (most of) these chapters helped expose weaknesses, areas in need of clarification, or counterarguments I hadn't considered. Foremost among these is Alan Jacobs Richardson (01d55), whose lengthy discussion with me regarding the relationship between Rebellion and E.T.A. Hoffmann's "The Nutcracker and the Mouse King" \cite{ref11} inspired significant chunks of Chapters 20 and 21.